% Mock midterm answer walkthrough.
% Compile from the repository root:
%   pdflatex -interaction=nonstopmode -halt-on-error -output-directory=output/pdf Lecture8/mock-mid-term-answers.tex
% This file also compiles directly from inside Lecture8/.
\documentclass[aspectratio=169]{beamer}

\usepackage{amsmath,amssymb,mathtools}
\usepackage{array,booktabs}
\usepackage{xcolor}
\usepackage{iftex}
\ifPDFTeX
  \usepackage[UTF8,heading=false]{ctex}
\else
  \usepackage{fontspec}
  \usepackage{xeCJK}
  \setCJKmainfont{Microsoft YaHei}
\fi

% The course theme is at the repository root.
\makeatletter
\def\input@path{{./}{../}{Lecture8/}}
\makeatother

\mode<presentation>{
  \usetheme{OSUCSE169}
  \setbeamercovered{invisible}
}
\setbeamertemplate{navigation symbols}{}
\setbeamertemplate{footline}[frame number]
\setbeamercolor{headline}{fg=black,bg=black}
\setbeamercolor{section in head/foot}{bg=gray}
\setbeamertemplate{headline}{%
  \begin{beamercolorbox}[ht=2.25ex,dp=3.75ex]{section in head/foot}
    \insertnavigation{\paperwidth}
  \end{beamercolorbox}}

\definecolor{lecturegreen}{RGB}{48,105,21}
\definecolor{lectureblue}{RGB}{0,76,170}
\definecolor{lectureorange}{RGB}{202,105,0}

\newcommand{\R}{\mathbb{R}}
\newcommand{\vect}[1]{\mathbf{#1}}
\newcommand{\Span}{\operatorname{Span}}
\newcommand{\Col}{\operatorname{Col}}
\newcommand{\rank}{\operatorname{rank}}
\newcommand{\sectioncard}[2]{%
  \begin{frame}[plain]
    \vfill\centering
    {\Huge\color{lectureorange}\bfseries #1}\par\bigskip
    {\LARGE\color{lecturegreen}\bfseries #2}
    \vfill
  \end{frame}}
\newcommand{\answerbox}[1]{%
  \begin{beamercolorbox}[sep=1.0ex,wd=\linewidth]{block body}
    \textbf{Answer:} #1
  \end{beamercolorbox}}

\title{MATH 100: Linear Algebra}
\subtitle{Mock Midterm - Answers and Walkthrough}
\author{Yuqing Yang}
\institute{Macau University of Science and Technology}
\date{10/21/2026}

\begin{document}

\begin{frame}[plain]
  \titlepage
\end{frame}

\begin{frame}{Walkthrough goals}
For each question, identify
\begin{itemize}
  \item the idea being tested;
  \item the shortest reliable method;
  \item the steps that earn partial credit;
  \item the final answer in a form that can be checked quickly.
\end{itemize}
\end{frame}

\section{Multiple choice}
\sectioncard{Section I}{Multiple choice - 20 points}

\begin{frame}{MC 1 - Number of solutions}
\[
\left[\begin{array}{ccc|c}
1&0&2&3\\
0&1&-1&4\\
0&0&0&0
\end{array}\right]
\]
There is no contradictory row. The third column has no pivot, so $x_3$ is free.
\[
x_1=3-2x_3,
\qquad
x_2=4+x_3.
\]
\answerbox{\textbf{C.} Infinitely many solutions.}
\end{frame}

\begin{frame}{MC 2 - Guaranteed dependence}
Any set containing more than $n$ vectors in $\R^n$ is linearly dependent.
\[
4>3.
\]
Therefore four vectors in $\R^3$ cannot all be linearly independent.
\answerbox{\textbf{C.} Four vectors in $\R^3$.}
\end{frame}

\begin{frame}{MC 3 - Product dimensions}
\[
A_{3\times2}B_{2\times4}.
\]
The inner dimensions match, so the product exists. The outer dimensions give its size:
\[
(3\times2)(2\times4)\longrightarrow 3\times4.
\]
\answerbox{\textbf{D.} $AB$ is $3\times4$.}
\end{frame}

\begin{frame}{MC 4 - Inverse of a product}
The factors reverse order:
\[
(AB)^{-1}=B^{-1}A^{-1}.
\]
Check:
\[
(AB)(B^{-1}A^{-1})
=A(BB^{-1})A^{-1}
=AA^{-1}=I.
\]
\answerbox{\textbf{B.} $(AB)^{-1}=B^{-1}A^{-1}$.}
\end{frame}

\begin{frame}{MC 5 - Onto transformations}
For $T(\vect x)=A\vect x$ to map onto $\R^m$, the equation
\[
A\vect x=\vect b
\]
must be consistent for every $\vect b\in\R^m$.
This occurs exactly when $A$ has a pivot in every row.
\answerbox{\textbf{A.} $A$ has a pivot in every row.}
\end{frame}

\begin{frame}{MC 6 - Determinant changes}
Start with $\det A=2$.
\begin{align*}
\text{multiply one row by }-3 &\Longrightarrow \text{factor }-3,\\
\text{interchange two rows} &\Longrightarrow \text{factor }-1.
\end{align*}
Hence
\[
\det B=(-3)(-1)(2)=6.
\]
\answerbox{\textbf{B.} $\det B=6$.}
\end{frame}

\begin{frame}{MC 7 - Singular matrices}
If $A\vect x=\vect0$ has a nontrivial solution, then
\begin{itemize}
  \item the columns of $A$ are linearly dependent;
  \item $A$ does not have a pivot in every column;
  \item $A$ is not invertible.
\end{itemize}
\answerbox{\textbf{C.} $A\vect x=\vect0$ has a nontrivial solution.}
\end{frame}

\begin{frame}{MC 8 - Basis for the column space}
The pivot columns are columns 1 and 3.
\medskip

Use the corresponding columns from the \textbf{original} matrix, not the reduced matrix.
\[
\{\vect a_1,\vect a_3\}
\]
is linearly independent and spans $\Col A$.
\answerbox{\textbf{B.} $\{\vect a_1,\vect a_3\}$.}
\end{frame}

\begin{frame}{MC 9 - Number of free variables}
Rank $3$ means that $A$ has $3$ pivot columns. Since $A$ has $7$ columns,
\[
\#\text{ free variables}
=\#\text{ columns}-\#\text{ pivot columns}
=7-3=4.
\]
\answerbox{\textbf{C.} The system $A\vect x=\vect0$ has $4$ free variables.}
\end{frame}

\begin{frame}{MC 10 - A linear transformation}
\[
T(x,y)=(x+2y,3x-y)
=
\begin{bmatrix}
1&2\\3&-1
\end{bmatrix}
\begin{bmatrix}x\\y\end{bmatrix}.
\]
This is a matrix transformation, so it is linear. The other options contain a constant, a square, or a product of variables.
\answerbox{\textbf{C.} $T(x,y)=(x+2y,3x-y)$.}
\end{frame}

\section{Brief answers}
\sectioncard{Section II}{Brief answers - 10 points}

\begin{frame}{Brief 1 - Consistency with a parameter}
The last row represents
\[
0=h-5.
\]
The system is consistent exactly when the right side is zero:
\[
h-5=0.
\]
\answerbox{$\boxed{h=5}$.}
\end{frame}

\begin{frame}{Brief 2 - Pivot and free variables}
\[
\begin{bmatrix}
\boxed{1}&2&0&-1&0&3\\
0&0&\boxed{1}&4&0&-2\\
0&0&0&0&\boxed{1}&5\\
0&0&0&0&0&0
\end{bmatrix}
\]
Pivot columns: $1,3,5$. Free-variable columns: $2,4,6$.
\answerbox{$\boxed{3}$ pivot variables and $\boxed{3}$ free variables.}
\end{frame}

\begin{frame}{Brief 3 - Checking an inverse}
\[
A=\begin{bmatrix}1&2\\3&7\end{bmatrix},
\qquad
C=\begin{bmatrix}7&-2\\-3&1\end{bmatrix}.
\]
\[
AC=
\begin{bmatrix}
7-6&-2+2\\
21-21&-6+7
\end{bmatrix}
=I_2.
\]
Also $CA=I_2$.
\answerbox{Yes. $C=A^{-1}$.}
\end{frame}

\begin{frame}{Brief 4 - Membership in a span}
Solve
\[
c_1\vect v_1+c_2\vect v_2+c_3\vect v_3=\vect b.
\]
The coordinate equations are
\[
c_1+c_3=2,
\qquad
c_2+c_3=3,
\qquad
c_1+c_2=1.
\]
They give $(c_1,c_2,c_3)=(0,1,2)$.
\answerbox{Yes. $\vect b=\vect v_2+2\vect v_3$.}
\end{frame}

\begin{frame}{Brief 5 - Isolating a matrix factor}
Starting with
\[
ABC=P,
\]
multiply by $A^{-1}$ on the left:
\[
BC=A^{-1}P.
\]
Then multiply by $C^{-1}$ on the right:
\[
B=A^{-1}PC^{-1}.
\]
\answerbox{$\boxed{B=A^{-1}PC^{-1}}$.}
\end{frame}

\section{Exercises}
\sectioncard{Section III}{Exercises - 70 points}

\begin{frame}{Exercise 1 - Solving a linear system}
\small
\[
\left[\begin{array}{ccc|c}
1&2&-1&3\\
2&5&1&8\\
-1&-1&2&-1
\end{array}\right]
\xrightarrow[,R_3\leftarrow R_3+R_1,]{,R_2\leftarrow R_2-2R_1,}
\left[\begin{array}{ccc|c}
1&2&-1&3\\
0&1&3&2\\
0&1&1&2
\end{array}\right]
\]
\[
\xrightarrow{R_3\leftarrow R_3-R_2}
\left[\begin{array}{ccc|c}
1&2&-1&3\\
0&1&3&2\\
0&0&-2&0
\end{array}\right].
\]
Thus $z=0$, $y=2$, and $x=-1$.
\answerbox{$\boxed{(x,y,z)=(-1,2,0)}$.}
\end{frame}

\begin{frame}{Exercise 2 - Matrix multiplication}
\small
\[
AB=
\begin{bmatrix}1&-2&3\\0&4&-1\end{bmatrix}
\begin{bmatrix}2&1\\-1&3\\4&0\end{bmatrix}.
\]
\[
=\begin{bmatrix}
1(2)+(-2)(-1)+3(4)&1(1)+(-2)(3)+3(0)\\
0(2)+4(-1)+(-1)(4)&0(1)+4(3)+(-1)(0)
\end{bmatrix}
\]
\answerbox{$\displaystyle \boxed{AB=\begin{bmatrix}16&-5\\-8&12\end{bmatrix}}$.}
\end{frame}

\begin{frame}{Exercise 3 - Determinant}
\[
D=\begin{bmatrix}
2&0&0&0\\
1&3&4&0\\
0&2&-1&5\\
0&0&0&4
\end{bmatrix}.
\]
Expand along row 1, then along the last column of the remaining matrix:
\begin{align*}
\det D
&=2\det\begin{bmatrix}3&4&0\\2&-1&5\\0&0&4\end{bmatrix}\\
&=2(4)\det\begin{bmatrix}3&4\\2&-1\end{bmatrix}
=8(-3-8).
\end{align*}
\answerbox{$\boxed{\det D=-88}$.}
\end{frame}

\begin{frame}[shrink=4]{Exercise 4 - Finding an inverse}
\small
\[
\left[\begin{array}{ccc|ccc}
1&1&0&1&0&0\\
0&1&1&0&1&0\\
1&0&1&0&0&1
\end{array}\right]
\xrightarrow{R_3\leftarrow R_3-R_1}
\left[\begin{array}{ccc|ccc}
1&1&0&1&0&0\\
0&1&1&0&1&0\\
0&-1&1&-1&0&1
\end{array}\right]
\]
\[
\xrightarrow{R_3\leftarrow R_3+R_2}
\left[\begin{array}{ccc|ccc}
1&1&0&1&0&0\\
0&1&1&0&1&0\\
0&0&2&-1&1&1
\end{array}\right]
\longrightarrow
\left[\begin{array}{ccc|ccc}
1&0&0&\frac12&-\frac12&\frac12\\
0&1&0&\frac12&\frac12&-\frac12\\
0&0&1&-\frac12&\frac12&\frac12
\end{array}\right].
\]
\answerbox{$\displaystyle \boxed{A^{-1}=\frac12\begin{bmatrix}1&-1&1\\1&1&-1\\-1&1&1\end{bmatrix}}$.}
\end{frame}

\begin{frame}{Exercise 5 - LU factorization}
Use the elimination multipliers
\[
m_{21}=2,
\qquad
m_{31}=-1,
\qquad
m_{32}=1.
\]
The row replacements produce
\[
U=\begin{bmatrix}
2&4&-2\\
0&1&1\\
0&0&4
\end{bmatrix}.
\]
Place the multipliers below the diagonal of $L$:
\answerbox{$\displaystyle
\boxed{L=\begin{bmatrix}1&0&0\\2&1&0\\-1&1&1\end{bmatrix},
\quad
U=\begin{bmatrix}2&4&-2\\0&1&1\\0&0&4\end{bmatrix}}$.}
\end{frame}

\begin{frame}{Exercise 6 - Recovering the original determinant}
The three operations have determinant factors
\[
-1,
\qquad
3,
\qquad
1.
\]
Therefore
\[
\det B=(-1)(3)(1)\det A=-3\det A.
\]
Given $\det B=36$,
\[
36=-3\det A.
\]
\answerbox{$\boxed{\det A=-12}$.}
\end{frame}

\begin{frame}{Exercise 7(a) - Unique solutions}
The coefficient matrix is triangular:
\[
A(k)=\begin{bmatrix}
1&1&0\\
0&k-1&1\\
0&0&k+2
\end{bmatrix}.
\]
Hence
\[
\det A(k)=1(k-1)(k+2).
\]
The system has a unique solution exactly when the determinant is nonzero.
\answerbox{$\boxed{k\ne1,-2}$.}
\end{frame}

\begin{frame}{Exercise 7(b) - The case \texorpdfstring{$k=1$}{k=1}}
At $k=1$,
\[
\begin{cases}
x_1+x_2=2,\\
x_3=1,\\
3x_3=3.
\end{cases}
\]
The third equation repeats the second. Let $x_2=t$:
\[
x_1=2-t,
\qquad
x_3=1.
\]
\answerbox{Infinitely many solutions:
$\displaystyle \boxed{\vect x=\begin{bmatrix}2\\0\\1\end{bmatrix}
+t\begin{bmatrix}-1\\1\\0\end{bmatrix},\ t\in\R}$.}
\end{frame}

\begin{frame}{Exercise 7(c) - The case \texorpdfstring{$k=-2$}{k=-2}}
At $k=-2$, the final row of the augmented system is
\[
\left[\begin{array}{ccc|c}
0&0&0&3
\end{array}\right].
\]
This row represents the contradiction
\[
0=3.
\]
Equivalently,
\[
\rank A<\rank[A\mid\vect b].
\]
\answerbox{The system is inconsistent and has no solution.}
\end{frame}

\begin{frame}{Compact answer key}
\small
\textbf{Multiple choice}
\[
\boxed{\text{C, C, D, B, A, B, C, B, C, C}}
\]
\textbf{Brief answers}
\begin{enumerate}
  \item $h=5$.
  \item 3 pivot variables and 3 free variables.
  \item Yes, $C=A^{-1}$.
  \item $\vect b=\vect v_2+2\vect v_3$.
  \item $B=A^{-1}PC^{-1}$.
\end{enumerate}
\textbf{Exercises}
\[
(-1,2,0),\quad
\begin{bmatrix}16&-5\\-8&12\end{bmatrix},\quad
-88,\quad
A^{-1},\quad
A=LU,\quad
-12,
\]
with the parameter cases $k\ne1,-2$, $k=1$, and $k=-2$ classified above.
\end{frame}

\end{document}
